\section*{Open Questions}

The following five questions emerged from this replication as genuinely open
directions that build on --- rather than merely restate --- the paper's
contributions. Each has a concrete, free-endpoint next-step probe.

\begin{enumerate}

\item \textbf{Compact 2e ansatz vs modern adaptive ansatze on non-trivial 2e
systems.} Does the ansatz retain its advantage against ADAPT-VQE / iterative-QCC
/ qubit-ADAPT on H$_2$/HeH$^+$ in larger bases (6-31G, cc-pVDZ) where adaptive
methods can grow their operator pool? The paper predates the adaptive-ansatz
literature. \emph{Probe:} extend our PySCF+OpenFermion pipeline with an ADAPT
loop; score (energy, params, CNOTs) at chemical-accuracy convergence across
$\{$H$_2$, HeH$^+$, LiH$\} \times \{$STO-3G, 6-31G, cc-pVDZ$\}$; free Qiskit
statevector only.

\item \textbf{Noise-aware ansatz choice under a calibrated device noise model.}
Does the 6-CNOT vs 8-CNOT vs 14-CNOT depth gap actually translate to the
fidelity gap the paper implies under realistic depolarising / amplitude-damping
/ readout error at NISQ rates, with and without ZNE and symmetry verification?
Central to C4+C5 which we could not test on retired ibm-5/ibm-14.
\emph{Probe:} Qiskit Aer noise model calibrated from archived ibm-5/ibm-14
calibration snapshots; sweep error rate $p\in[10^{-4},10^{-2}]\times$ (three
ansatze) $\times$ (three mitigation policies); free local Aer; deliverable
$3{\times}3$ accuracy-floor heatmap.

\item \textbf{Basis transferability across bond-length regimes.} If natural
orbitals are frozen at one geometry $R_0$ and only $\theta$ is re-optimised
elsewhere, how far can $R$ drift before the frozen-basis compact ansatz misses
chemical accuracy? Directly relevant to reactive-dynamics use and active-space
embedding, unaddressed by the paper. \emph{Probe:} add \texttt{--nocap R0} to
\texttt{vqe\_h2\_compact.py}; sweep $R_0 \in \{0.735, 1.4, 2.0\}$~\AA{} and
compute the transferability window for each; free local PySCF; deliverable
transferability curve.

\item \textbf{Active-space embedding: does ansatz error survive orbital
selection?} When compact-2e is embedded via CASSCF(2,2) / DMRG-based orbital
selection / SHCI in molecules bigger than H$_2$, is the sub-picohartree
ansatz error hidden under CASSCF-projection error? \emph{Probe:} PySCF CASSCF(2,2)
on $\{$H$_2$O, LiH, BeH$_2\}$; decompose total error as [active-space projection]
+ [ansatz]; compare against full-molecule CCSD(T); free local; deliverable
error-decomposition table.

\item \textbf{ML-selected ansatz architecture.} Can a classifier trained on
$(\text{molecule}, \text{basis}, \text{geometry}, \text{target accuracy})$
predict the best of $\{$compact-2e, UCCSD, $k$-UpCCGSD, ADAPT-VQE, HEA$\}$,
and does compact-2e occupy a distinct applicability envelope (=exactly 2
electrons, small basis, near equilibrium) that ML can cleanly identify?
Informs whether the paper's contribution is niche vs a general primitive.
\emph{Probe:} build a $\sim$50-tuple dataset across
$\{$H$_2$, HeH$^+$, H$_3^+$, LiH, BeH$_2\}\times$
$\{$STO-3G, 6-31G$\}\times$ ansatze; train XGBoost / decision-tree classifier;
report feature importance and the resulting envelope; free sklearn local;
deliverable a decision-tree diagram.

\end{enumerate}
